\documentclass[10pt,a4paper]{amsart}
\usepackage{amsmath,amssymb,amsthm,graphicx}
\usepackage{verbatim}
\newenvironment{reference}{\comment}{\endcomment}
%\newenvironment{reference}{}{}
\newenvironment{slogan}{\comment}{\endcomment}
\newenvironment{history}{\comment}{\endcomment}
% For commutative diagrams we use Xy-pic
\usepackage[all]{xy}

% We use 2cell for 2-commutative diagrams.
\xyoption{2cell}
% Theorem environments.
%
\theoremstyle{plain}
\newtheorem{theorem}[subsection]{Theorem}
\newtheorem{proposition}[subsection]{Proposition}
\newtheorem{lemma}[subsection]{Lemma}

\theoremstyle{definition}
\newtheorem{definition}[subsection]{Definition}
\newtheorem{example}[subsection]{Example}
\newtheorem{exercise}[subsection]{Exercise}
\newtheorem{situation}[subsection]{Situation}

\theoremstyle{remark}
\newtheorem{remark}[subsection]{Remark}
\newtheorem{remarks}[subsection]{Remarks}

\numberwithin{equation}{subsection}

% Macros
%
\def\lim{\mathop{\mathrm{lim}}\nolimits}
\def\colim{\mathop{\mathrm{colim}}\nolimits}
\def\Spec{\mathop{\mathrm{Spec}}}
\def\Hom{\mathop{\mathrm{Hom}}\nolimits}
\def\Ext{\mathop{\mathrm{Ext}}\nolimits}
\def\SheafHom{\mathop{\mathcal{H}\!\mathit{om}}\nolimits}
\def\SheafExt{\mathop{\mathcal{E}\!\mathit{xt}}\nolimits}
\def\Sch{\mathit{Sch}}
\def\Mor{\mathop{\mathrm{Mor}}\nolimits}
\def\Ob{\mathop{\mathrm{Ob}}\nolimits}
\def\Sh{\mathop{\mathit{Sh}}\nolimits}
\def\NL{\mathop{N\!L}\nolimits}
\def\CH{\mathop{\mathrm{CH}}\nolimits}
\def\proetale{{pro\text{-}\acute{e}tale}}
\def\etale{{\acute{e}tale}}
\def\QCoh{\mathit{QCoh}}
\def\Ker{\mathop{\mathrm{Ker}}}
\def\Im{\mathop{\mathrm{Im}}}
\def\Coker{\mathop{\mathrm{Coker}}}
\def\Coim{\mathop{\mathrm{Coim}}}

% Boxtimes
%
\DeclareMathSymbol{\boxtimes}{\mathbin}{AMSa}{"02}

%
% Macros for moduli stacks/spaces
%
\def\QCohstack{\mathcal{QC}\!\mathit{oh}}
\def\Cohstack{\mathcal{C}\!\mathit{oh}}
\def\Spacesstack{\mathcal{S}\!\mathit{paces}}
\def\Quotfunctor{\mathrm{Quot}}
\def\Hilbfunctor{\mathrm{Hilb}}
\def\Curvesstack{\mathcal{C}\!\mathit{urves}}
\def\Polarizedstack{\mathcal{P}\!\mathit{olarized}}
\def\Complexesstack{\mathcal{C}\!\mathit{omplexes}}
% \Pic is the operator that assigns to X its picard group, usage \Pic(X)
% \Picardstack_{X/B} denotes the Picard stack of X over B
% \Picardfunctor_{X/B} denotes the Picard functor of X over B
\def\Pic{\mathop{\mathrm{Pic}}\nolimits}
\def\Picardstack{\mathcal{P}\!\mathit{ic}}
\def\Picardfunctor{\mathrm{Pic}}
\def\Deformationcategory{\mathcal{D}\!\mathit{ef}}
\usepackage{cleveref}
\crefname{originalreference}{original reference}{original references}
\Crefname{originalreference}{Original reference}{Original references}
\makeatletter
\newlabel{derived-lemma-two-ss-complex-functor}{{derived-lemma-two-ss-complex-functor}{source}{}{sourcepointer-1}{}}
\newlabel{derived-lemma-two-ss-complex-functor@cref}{{[originalreference][0][]derived-lemma-two-ss-complex-functor}{[1][0][]source}{}{}{}}
\makeatother
\crefname{originalreference}{original reference}{original references}
\Crefname{originalreference}{Original reference}{Original references}
\makeatletter
\newlabel{dualizing-proposition-structure-injectives-noetherian}{{dualizing-proposition-structure-injectives-noetherian}{source}{}{sourcepointer-2}{}}
\newlabel{dualizing-proposition-structure-injectives-noetherian@cref}{{[originalreference][0][]dualizing-proposition-structure-injectives-noetherian}{[1][0][]source}{}{}{}}
\makeatother
\title{The bounded below derived category of modules supported in a specialization-stable subset is equivalent to the full subcategory of complexes supported there}
\author{The Stacks Project Authors}
\date{stacks-project:local-cohomology.tex:local-snapshot-20261012; GNU FDL 1.2 or later}
\begin{document}
\maketitle
\noindent\textit{Editorial note.} All delivered mathematical text is the author's own, copied from the pinned Stacks Project chapter source sources/stacks-local-cohomology/pinned-original-source.tex; no statement or proof step was rewritten, reordered or removed.\par\medskip
\noindent\textit{Editorial note.} Preamble provenance: the Stacks Project keeps one single shared preamble for every chapter; this card takes it directly from the staged sources/stacks-local-cohomology/preamble.tex as three verbatim spans (lines 8-12, 14-18 and 154-217), so no derived file was written and nothing was added or redefined. The wrapper class and the base packages are supplied by the assembler.\par\medskip
\noindent\textit{Editorial note.} Bibliography provenance: neither the counted statement nor its proof contains a citation (this lane rescans the whole span before writing the card), so the card carries no bibliography and no entry was reconstructed.\par\medskip
\noindent\textit{Editorial note.} Selection provenance: this card was opened under the file-based control contract (CONTROL.md), with the mechanical pre-card gate run on its statement and proof interval before the card was written.\par\medskip
\noindent\textit{Allowed standard prerequisites (editorial).}\begin{itemize}
\item The structure of injective modules over a Noetherian ring as direct sums of injective hulls of residue fields. Invoked inside the counted proof at source lines 1028-1031 and attributed to a proposition of another chapter (a cross-chapter pointer). Identifies the zeroth local cohomology of an injective module and shows that it is again injective.
\item The spectral sequence of a composite functor in the form available for complexes. Invoked inside the counted proof at source lines 1043-1047 as a cross-chapter pointer. Shows that the unit of the adjunction is an isomorphism on bounded below complexes with support, given that only the q = 0 terms survive.
\end{itemize}\medskip
\ifdefined\hypertarget\hypertarget{sourcepointer-1}{}\fi
\noindent\textit{Source reference.} Original reference derived-lemma-two-ss-complex-functor: sources/stacks-local-cohomology/pinned-original-source.tex, source lines 1047--1047. Navigation pointer only: the counted text refers to the Stacks item derived-lemma-two-ss-complex-functor, which belongs to another chapter of the project and is not part of this card. The pointer carries the target's own label as its number because this chapter's staged snapshot does not contain the other chapter's tag table; its content is not reproduced here.\par\medskip
\ifdefined\hypertarget\hypertarget{sourcepointer-2}{}\fi
\noindent\textit{Source reference.} Original reference dualizing-proposition-structure-injectives-noetherian: sources/stacks-local-cohomology/pinned-original-source.tex, source lines 1030--1030. Navigation pointer only: the counted text refers to the Stacks item dualizing-proposition-structure-injectives-noetherian, which belongs to another chapter of the project and is not part of this card. The pointer carries the target's own label as its number because this chapter's staged snapshot does not contain the other chapter's tag table; its content is not reproduced here.\par\medskip
\begin{lemma}
\label{lemma-equal-plus}
Let $A$ be a Noetherian ring. Let $T \subset \Spec(A)$ be a subset stable
under specialization. The functor $D^+(\text{Mod}_{A, T}) \to D^+_T(A)$
is an equivalence.
\end{lemma}

\begin{proof}
Let $M$ be an object of $\text{Mod}_{A, T}$. Choose an embedding
$M \to J$ into an injective $A$-module. By
Dualizing Complexes, Proposition
\ref{dualizing-proposition-structure-injectives-noetherian}
the module $J$ is a direct sum of injective hulls of residue fields.
Let $E$ be an injective hull of the residue field of $\mathfrak p$.
Since $E$ is $\mathfrak p$-power torsion we see that
$H^0_T(E) = 0$ if $\mathfrak p \not \in T$ and
$H^0_T(E) = E$ if $\mathfrak p \in T$.
Thus $H^0_T(J)$ is injective as a direct sum of injective hulls
(by the proposition) and we have an embedding $M \to H^0_T(J)$.
Thus every object $M$ of $\text{Mod}_{A, T}$ has an injective resolution
$M \to J^\bullet$ with $J^n$ also in $\text{Mod}_{A, T}$. It follows
that $RH^0_T(M) = M$.

\medskip\noindent
Next, suppose that $K \in D_T^+(A)$. Then the spectral sequence
$$
R^qH^0_T(H^p(K)) \Rightarrow R^{p + q}H^0_T(K)
$$
(Derived Categories, Lemma \ref{derived-lemma-two-ss-complex-functor})
converges and above we have seen that only the terms with $q = 0$
are nonzero. Thus we see that $RH^0_T(K) \to K$ is an isomorphism.
Thus the functor $D^+(\text{Mod}_{A, T}) \to D^+_T(A)$
is an equivalence with quasi-inverse given by $RH^0_T$.
\end{proof}

\end{document}
